\section{Conclusion}\label{sec:conclusion}

If local summaries preserve the evidence a task oracle needs, and local
merges preserve the cross-span interactions, then preference
optimization can be carried out on compression trees rather than full
documents. The neural-operator layer now makes explicit where the
learned architecture enters: discretization invariance and compact-set
approximation for the equation-(6) operator family explain how an ideal
operator can be realized to a chosen tolerance, the transformer-inclusion
argument identifies finite attention stacks with that architecture
surface, and the realization bridge of
Proposition~\ref{prop:neural-operator-bridge} converts that tolerance
into approximate local-law budgets. The oracle/projection objective
supplies the training interpretation: fit the task oracle while
weighting C1/C2/C3 as a projection toward the theorem-eligible
subspace. The projection/structured-error equivalence says this is the
same as assuming the approximation error is structured so that its zero
set is exactly local-law validity. Everything downstream then reuses
the existing theorem-backed route: the structural bundle gives root
preservation, the optimization bundle gives objective equivalence for
DPO/GRPO, and the certificate bundle turns sampled audits into the
finite-sample gap decomposition laid out in Section~\ref{sec:discussion}:
an empirical IPW estimate plus a fixed calibration residual plus an
estimation envelope that shrinks monotonically in the total number of
labels collected, global or local. The four statements the Discussion
uses (the main gap bound, fixed calibration, envelope rate, and
global/local decomposition) are cross-referenced to their formalized
versions in Appendix~\ref{app:proof-artifacts}.

The argument arcs from the classical to the learned case along a single
gradient. Section~\ref{sec:mergeable-sketches} grounds the framework in
the data-systems literature on mergeable sketches, with HLL as the
empirical representative and a broader encode/merge/decode class
covering the rest.
Sections~\ref{sec:markov-interlude}--\ref{sec:markov-walkthrough} show
a task whose oracle defeats unordered register-wise merges without
boundary state and a tree-plus-FNO
instance of the more general neural-operator realization layer that
resolves it. Section~\ref{sec:hll-parity} closes the
classical-limit loop empirically: running the framework with a
deterministic register-wise merge recovers HyperLogLog byte-for-byte on
a register-based implementation, and a fully end-to-end learned version
tracks the classical sketch within a small multiple of its theoretical
relative standard error. Appendix~\ref{app:operator-overlap} records the
formal overlap: transformer stacks sit inside equation-(6) neural
operators, mergeable sketches sit inside exact local-law operators, and
the certified intersection is the mergeable neural-operator subspace.
Section~\ref{sec:social-measurement} sets up the social-science
measurement setting, and Section~\ref{sec:manifesto-llm} then takes the
same machinery to Manifesto-RILE scoring, where no hand-designed sketch
exists and the sufficient statistic is a high-dimensional learned
representation. The Manifesto section also adds the prompt-program case: the
summarizer and scorer can be updated repeatedly, but the empirical claim
is still local preservation, not blind faith in a full-context read.
Section~\ref{sec:manifesto-qsentence} completes the arc with genuine
node-level supervision: the corpus's quasi-sentence expert codes give a
gold oracle at every node, a trained neural-operator merge composes on
every signal-dense policy dimension, and the failures land exactly
where the label system predicts them (sparse dimensions, constructs the
codes do not determine, objectives that do not price the local laws).
The main finding: for
tasks whose structure aligns with the local laws, reliable measurement
on long documents becomes a question of local architecture and audit
budget rather than of full-document annotation alone. The trade
restricts the framework's scope to locally-decomposable oracles, and it
pays off in the many social-scientific settings where such
decomposition is natural.

The Semantic Forests companion extends these ideas to orchestration
and deployment at system scale.
